% Fragment: certified convex recourse, value-factor curvature, recognition of
% affine responses, and the negative-curvature target.
% No preamble.  Uses \R \Z \E \norm \dist \poly \OPT (= F^*) and the labels
% lem:round, prop:filter, lem:contraction, lem:localize, lem:states,
% thm:approx, thm:exact of the main text.

\section{Certified convex recourse}\label{sec:cr}

Large positive curvature often sits in variables that enter the objective
convexly once the remaining variables are fixed. This section eliminates such
variables exactly, as convex value factors, and shows that the
corrected-grid algorithm then depends on the upper coordinate curvature of
the \emph{reduced} objective. One curvature lemma covers three cases: the
direct bound, globally affine responses, and piecewise-affine responses.
The section ends with the limits of this route.

\subsection{Model}

\begin{definition}[Convex recourse model]\label{def:cr-model}
The variables split into retained variables $z\in Z=\prod_{i\in\mathcal R}Z_i$,
where each $Z_i$ is $[l_i,u_i]$ or $[l_i,u_i]\cap\Z$ as in the main model, and
private blocks $y_t\in Y_t$, $t=1,\ldots,T$, where
$Y_t=\{y\in\R^{n_t}:G_ty\le\gamma_t\}$ is a nonempty bounded rational polytope.
Private variables are continuous. Write $\widehat Z$ for the continuous hull of
$Z$. Each block attaches to an ordered set $S_t\subseteq\mathcal R$ of
$k_t$ retained coordinates; for $i\in S_t$ let $j(t,i)$ be its position in
$S_t$. The objective is
\begin{equation}\label{eq:cr-model}
F(z,y)=q_0(z)+\sum_{t=1}^T f_t(y_t,z_{S_t}),\qquad
f_t(y,v)=\tfrac12y^{\mathsf T}C_ty+y^{\mathsf T}K_tv+c_t^{\mathsf T}y,
\end{equation}
with rational data, $C_t\succeq0$, and a rational quadratic
$q_0(z)=\sum_{a\in\mathcal A}f_a(z_{S_a})$ with Hessian $H_0$. All terms that
involve retained variables only are placed in $q_0$; thus $(H_0)_{ii}$ is the
diagonal entry of the full Hessian of $F$ at $z_i$. Fixed coordinates are
substituted first, so every interval $[l_i,u_i]$ has positive length.
A tree decomposition of the retained variables is supplied, with maximum bag
size $p$ and $N$ nodes, such that every $S_a$ and every $S_t$ lies in a bag.
\end{definition}

The full Hessian of $F$ need not be positive semidefinite, and $C_t$ may be
singular. Distinct private blocks interact only through $z$. Define the
\emph{value factors} and the reduced objective
\begin{equation}\label{eq:cr-value}
\phi_t(v)=\min_{y\in Y_t}f_t(y,v)\quad(v\in\R^{k_t}),\qquad
V(z)=q_0(z)+\sum_t\phi_t(z_{S_t}).
\end{equation}
Since the blocks are independent once $z$ is fixed,
$V(z)=\min_yF(z,y)$ for every $z\in\widehat Z$, and $\min_ZV=\OPT:=\min F$.
Every $z$ lifts to $(z,y(z))$ with $y_t(z)\in\arg\min_{Y_t}f_t(\cdot,z_{S_t})$
and $F(z,y(z))=V(z)$.

\subsection{Leaf certificates and the curvature of a value factor}

Fix one block and drop the index $t$: $y\in\R^n$, $v\in\R^k$,
$f(y,v)=\frac12y^{\mathsf T}Cy+y^{\mathsf T}Kv+c^{\mathsf T}y$,
$Y=\{y:Gy\le\gamma\}$ with $m$ rows, and $\phi(v)=\min_{y\in Y}f(y,v)$.

\begin{definition}[Leaf certificate]\label{def:leaf}
Let $P\subseteq\R^k$ be a closed polytope with nonempty interior. A
\emph{certificate on $P$} consists of affine maps $\bar y(v)=a+Bv$ and
$\lambda(v)=\lambda_0+\Lambda v\in\R^m$ such that, for all $v\in P$,
\begin{align}
&G\bar y(v)\le\gamma,\qquad \lambda(v)\ge0,\label{eq:leaf-feas}\\
&C\bar y(v)+Kv+c+G^{\mathsf T}\lambda(v)=0,\label{eq:leaf-stat}\\
&\lambda_k(v)\bigl(\gamma_k-G_k\bar y(v)\bigr)=0\qquad(k=1,\ldots,m).\label{eq:leaf-comp}
\end{align}
A \emph{certificate on a box $\Pi\subseteq\R^k$} is a finite family of closed
polytopes $P_1,\ldots,P_R\subseteq\Pi$ with nonempty interiors and union $\Pi$,
each with a certificate. For a private box $Y=[l,u]$ the rows are
$y\le u$ and $-y\le-l$, and $\lambda$ splits into upper and lower bound
multipliers.
\end{definition}

A certificate is checked exactly in time polynomial in its encoding length.
Stationarity~\eqref{eq:leaf-stat} is an identity between affine maps.
Complementarity~\eqref{eq:leaf-comp} is a quadratic polynomial vanishing on a
set with nonempty interior, which holds if and only if all its coefficients
vanish. An affine inequality $\alpha+\beta^{\mathsf T}v\ge0$ on a nonempty
polytope $\{v:Wv\le\omega\}$ holds if and only if some $\pi\ge0$ satisfies
$W^{\mathsf T}\pi=-\beta$ and $\omega^{\mathsf T}\pi\le\alpha$: sufficiency is
$\alpha+\beta^{\mathsf T}v=\alpha-\pi^{\mathsf T}Wv\ge\alpha-\pi^{\mathsf T}\omega\ge0$,
and necessity is linear programming duality for
$\min\{\beta^{\mathsf T}v:Wv\le\omega\}$. Such rational multipliers are part of
the certificate. When the leaves are produced by a binary space partition
(each internal node splits its polytope by one rational hyperplane into two
closed halves with nonempty interiors), induction on the tree shows that the
leaves cover the root. On a box, the range of $\alpha+\beta^{\mathsf T}v$ is
also available in closed form from the signs of $\beta$.

\begin{lemma}[Certificate identity and piece Hessian]\label{lem:leaf}
Let $(\bar y,\lambda)$ be a certificate on $P$. Then:
\begin{enumerate}
\item[(a)] For $v\in P$ and $y\in\R^{n}$,
\begin{equation}\label{eq:kkt-identity}
f(y,v)-f(\bar y(v),v)=\tfrac12(y-\bar y(v))^{\mathsf T}C(y-\bar y(v))
+\lambda(v)^{\mathsf T}(\gamma-Gy).
\end{equation}
The right side is nonnegative for $y\in Y$. Hence $\phi=q_P$ on $P$, where
$q_P(v)=f(\bar y(v),v)$.
\item[(b)] Let $\mathcal K=\{k:\lambda_k\not\equiv0\}$. Then $G_kB=0$ for
$k\in\mathcal K$, and
\[
\nabla^2q_P=B^{\mathsf T}CB+B^{\mathsf T}K+K^{\mathsf T}B=-B^{\mathsf T}CB\preceq0.
\]
\item[(c)] Put $\mathcal T=\{w:G_kw=0\ (k\in\mathcal K)\}$. For every
coordinate $j$,
\begin{equation}\label{eq:energy}
(B^{\mathsf T}CB)_{jj}
=-\min_{w\in\mathcal T}\bigl(w^{\mathsf T}Cw+2w^{\mathsf T}Ke_j\bigr).
\end{equation}
\end{enumerate}
\end{lemma}

\begin{proof}
(a) The function $f(\cdot,v)$ is quadratic with Hessian $C$, and
by~\eqref{eq:leaf-stat} its gradient at $\bar y=\bar y(v)$ is
$-G^{\mathsf T}\lambda$. Exact expansion gives
\[
f(y,v)-f(\bar y,v)=\lambda^{\mathsf T}(G\bar y-Gy)+\tfrac12(y-\bar y)^{\mathsf T}C(y-\bar y),
\]
and $\lambda^{\mathsf T}(G\bar y-Gy)=\lambda^{\mathsf T}(\gamma-Gy)-\lambda^{\mathsf T}(\gamma-G\bar y)
=\lambda^{\mathsf T}(\gamma-Gy)$ by~\eqref{eq:leaf-comp}. For $y\in Y$ both
terms on the right of~\eqref{eq:kkt-identity} are nonnegative, since
$C\succeq0$ and $\lambda(v)\ge0$. Equality holds at $y=\bar y(v)\in Y$.

(b) For each $k$, the polynomial $\lambda_k(v)(\gamma_k-G_k a-G_kBv)$ vanishes
on $P$, which has nonempty interior, so it is the zero polynomial. The
polynomial ring is an integral domain; for $k\in\mathcal K$ the first factor
is a nonzero polynomial, so the second is zero and $G_kB=0$. The $v$-linear
part of~\eqref{eq:leaf-stat} reads $CB+K+G^{\mathsf T}\Lambda=0$, and row $k$
of $\Lambda$ vanishes for $k\notin\mathcal K$. Hence
\[
B^{\mathsf T}(CB+K)=-\sum_{k\in\mathcal K}(G_kB)^{\mathsf T}\Lambda_k=0,
\]
so $B^{\mathsf T}K=-B^{\mathsf T}CB$, which is symmetric; therefore also
$K^{\mathsf T}B=-B^{\mathsf T}CB$. Substituting $y=a+Bv$ in $f$ gives
$\nabla^2q_P=B^{\mathsf T}CB+B^{\mathsf T}K+K^{\mathsf T}B=-B^{\mathsf T}CB$.

(c) The function $\Phi(w)=w^{\mathsf T}Cw+2w^{\mathsf T}Ke_j$ is convex. By (b),
$Be_j\in\mathcal T$, and for $w\in\mathcal T$,
\[
\tfrac12\nabla\Phi(Be_j)^{\mathsf T}w=(CBe_j+Ke_j)^{\mathsf T}w
=-\sum_{k\in\mathcal K}(\Lambda_ke_j)(G_kw)=0.
\]
A convex function whose gradient at a point of a linear subspace is
orthogonal to that subspace attains its minimum over the subspace there. Thus
$\min_{\mathcal T}\Phi=\Phi(Be_j)=(B^{\mathsf T}CB)_{jj}+2(B^{\mathsf T}K)_{jj}
=-(B^{\mathsf T}CB)_{jj}$.
\end{proof}

Formula~\eqref{eq:energy} says that a piece cancels direct curvature exactly
by the energy that the private response spends to follow a unit change of
$v_j$, using only directions that keep every constraint with a nonzero
multiplier active. If all private coordinates coupled to $v_j$ are held at
bounds with nonzero multipliers, the piece cancels nothing.

\begin{proposition}[Coordinate curvature of a value factor]\label{prop:vf-curv}
\begin{enumerate}
\item[(a)] $\phi$ is finite and concave on $\R^k$.
\item[(b)] Let $(P_r,\bar y_r=a_r+B_rv,\lambda_r)_{r=1}^R$ be a certificate on a
box $\Pi$, and put
\[
\sigma_j=\min_r\,(B_r^{\mathsf T}CB_r)_{jj}\ \ge0.
\]
For every $v\in\Pi$ and every interval $I$ with $v+Ie_j\subseteq\Pi$, the
function $t\mapsto\phi(v+te_j)+\sigma_jt^2/2$ is concave on $I$.
\item[(c)] If a number $\sigma'$ has the property in (b), then
$\sigma'\le\sigma_j$. In particular $\sigma_j$ is the same for every
certificate on $\Pi$.
\end{enumerate}
\end{proposition}

\begin{proof}
(a) For fixed $y$, $f(y,v)$ is affine in $v$. Thus $\phi$ is a pointwise
minimum of affine functions over the compact nonempty set $Y$, hence finite
and concave.

(b) Let $I=[\alpha,\beta]$ with $\alpha<\beta$; a point interval is trivial.
For each $r$, $J_r=\{t\in I:v+te_j\in P_r\}$ is a closed interval, possibly
empty or a point, and $\bigcup_rJ_r=I$. The union of the nondegenerate $J_r$
is closed, and its complement in $I$ is relatively open and contained in the
finite union of the degenerate $J_r$. A nonempty relatively open subset of a
nondegenerate interval is infinite, so the nondegenerate $J_r$ cover $I$. Let
$\alpha=t_0<\cdots<t_s=\beta$ list $\alpha$, $\beta$ and all endpoints of
nondegenerate $J_r$. The midpoint of $[t_{q-1},t_q]$ lies in some nondegenerate
$J_r$; because no endpoint of $J_r$ lies strictly between $t_{q-1}$ and $t_q$,
we get $[t_{q-1},t_q]\subseteq J_r$. By Lemma~\ref{lem:leaf}, on this subinterval
$\phi(v+te_j)=q_{P_r}(v+te_j)$, a quadratic in $t$ with second derivative
$-(B_r^{\mathsf T}CB_r)_{jj}\le-\sigma_j$. Hence
$\psi(t)=\phi(v+te_j)+\sigma_jt^2/2$ is concave on each $[t_{q-1},t_q]$. It is
continuous on $I$, because the concave function $\phi$ is continuous on $\R^k$.
At an interior breakpoint, $\psi'_+(t_q)-\psi'_-(t_q)$ equals the jump of the
one-sided derivatives of the concave function $t\mapsto\phi(v+te_j)$, which is
nonpositive; the smooth term $\sigma_jt^2/2$ has no jump. Therefore the right
derivative of $\psi$ is nonincreasing on $[\alpha,\beta)$. Since $\psi$ is
continuous and piecewise $C^1$, $\psi(t)=\psi(\alpha)+\int_\alpha^t\psi'_+$, so
$\psi$ is concave.

(c) Take any $r$ and an interior point $v_0$ of $P_r$. For small $|t|$,
$\phi(v_0+te_j)=q_{P_r}(v_0+te_j)$, whose second derivative is
$-(B_r^{\mathsf T}CB_r)_{jj}$. Concavity of $\phi(v_0+te_j)+\sigma't^2/2$ forces
$\sigma'\le(B_r^{\mathsf T}CB_r)_{jj}$. Minimizing over $r$ gives
$\sigma'\le\sigma_j$. Applying this to two certificates shows that each
$\sigma_j$ is at most the other.
\end{proof}

Part~(a) alone certifies $\sigma_j=0$ for every block; this is the direct
bound. Part~(c) shows that a piecewise certificate loses nothing: the
cancellation it certifies is the largest one possible for coordinatewise
semiconcavity of this factor.

\begin{proposition}[Finite certificates exist]\label{prop:cert-exist}
Let $Y=[l,u]$ with $l<u$, and let $\Pi\subseteq\R^k$ be a rational box with
nonempty interior. Then $\phi$ has a certificate on $\Pi$ whose polytopes are
the leaves of a rational binary space partition.
\end{proposition}

\begin{proof}
For a pattern $\sigma\in\{\mathrm L,\mathrm F,\mathrm U\}^n$ let $\Omega_\sigma$
be the set of $(v,y,\lambda^{\mathrm{low}},\lambda^{\mathrm{up}})$ with $v\in\Pi$,
$l\le y\le u$, $\lambda^{\mathrm{low}},\lambda^{\mathrm{up}}\ge0$,
$Cy+Kv+c=\lambda^{\mathrm{low}}-\lambda^{\mathrm{up}}$, and, coordinatewise,
$y_i=l_i,\lambda^{\mathrm{up}}_i=0$ if $\sigma_i=\mathrm L$;
$\lambda^{\mathrm{low}}_i=\lambda^{\mathrm{up}}_i=0$ if $\sigma_i=\mathrm F$;
$y_i=u_i,\lambda^{\mathrm{low}}_i=0$ if $\sigma_i=\mathrm U$. Each $\Omega_\sigma$
is a rational polyhedron, and every point of it satisfies the KKT conditions,
including complementarity. Its projection $Q_\sigma$ to the $v$ coordinates
is a rational polytope. Every $v\in\Pi$ has an optimal $y$ and, since the
constraints are linear, KKT multipliers; these fit some pattern. Thus the
$Q_\sigma$ cover $\Pi$. The union of the full-dimensional $Q_\sigma$ is closed
and contains $\Pi$ minus finitely many lower-dimensional polytopes, a dense
subset of $\Pi$; hence it equals $\Pi$.

Triangulate each full-dimensional $Q_\sigma$ into $k$-simplices with rational
vertices. At each vertex $w$ of a simplex $\Delta$, choose a rational point
$(w,y_w,\lambda_w)\in\Omega_\sigma$. On $\Delta$ define $(\bar y,\lambda)$ by
affine interpolation of these vertex values in barycentric coordinates. For
$v\in\Delta$ the point $(v,\bar y(v),\lambda(v))$ is a convex combination of
points of the convex set $\Omega_\sigma$, so it lies in $\Omega_\sigma$. Hence
$(\bar y,\lambda)$ is a certificate on $\Delta$.

Let $\mathcal H$ be the finite set of hyperplanes spanned by facets of all
these simplices. Starting from $\Pi$, split every current leaf whose interior
meets a hyperplane of $\mathcal H$ by that hyperplane, until no leaf interior
meets any of them. Both halves of every split have nonempty interior. Let $L$
be a final leaf. Its interior is connected and misses every simplex boundary,
so for each simplex $\Delta$ either $\operatorname{int}L\subseteq\operatorname{int}\Delta$
or $\operatorname{int}L\cap\Delta=\emptyset$. The full-dimensional simplices
cover $\Pi$, so the first case occurs for some $\Delta$, and $L\subseteq\Delta$.
Restricting the certificate of $\Delta$ to $L$ completes the proof.
\end{proof}

The same argument applies to a polytope $Y$, with patterns given by subsets
of active rows. The construction can be exponential; no bound on the number
of leaves is claimed. By Proposition~\ref{prop:vf-curv}(c), however, the
best cancellation is always certifiable by some finite certificate.

\subsection{The recourse theorem}

\begin{lemma}[Growth transfer]\label{lem:cr-growth}
Suppose $F(z,y)-\OPT\ge g\bigl(\norm{z-z^*}^2+\sum_t\norm{y_t-y_t^*}^2\bigr)$ on
the feasible set, with $g>0$. Then $(z^*,y^*)$ is the unique minimizer, $z^*$
is the unique minimizer of $V$, and for every $z\in Z$ and every choice of
optimal responses $y(z)$,
\[
V(z)-\OPT\ge g\Bigl(\norm{z-z^*}^2+\sum_t\norm{y_t(z)-y_t^*}^2\Bigr)
\ge g\norm{z-z^*}^2.
\]
If block $t$ has a single certificate $\bar y_t=a_t+B_tv$ on the whole
attachment box, then $y_t^*=\bar y_t(z^*_{S_t})$, and summing over such blocks
gives the metric $I+\sum_tE_t^{\mathsf T}B_t^{\mathsf T}B_tE_t\succeq I$, where
$E_tz=z_{S_t}$.
\end{lemma}

\begin{proof}
Growth with $g>0$ makes the minimizer unique. Evaluate it at $(z,y(z))$ and use
$F(z,y(z))=V(z)$. A minimizer $z'$ of $V$ lifts to a minimizer of $F$, so
$z'=z^*$. With an affine certificate, $(z^*,\bar y(z^*_{S}))$ has value
$V(z^*)=\OPT$, so uniqueness gives $y_t^*=\bar y_t(z^*_{S_t})$ and
$\bar y_t(z_{S_t})-y_t^*=B_tE_t(z-z^*)$.
\end{proof}

The following two lemmas replace the quadratic-denominator argument in the
proof of Theorem~\ref{thm:approx} and the box height bound of
Lemma~\ref{lem:height}.

\begin{lemma}[Exact tables with value factors]\label{lem:cr-bits}
Let every grid coordinate at a stage have bit length at most $b$. Then each
value $\phi_t(z_{S_t})$ at a grid assignment is a rational of bit length
$\poly(S+b)$ and is computable in that time, where $S$ is the encoding length
of the data. Every message entry, calibrated table entry and coordinate
min-marginal is a sum of at most $m_0=|\mathcal A|+T+|\mathcal R|$ such values
and corrections, hence has bit length at most $m_0\poly(S+b)$.
\end{lemma}

\begin{proof}
A convex quadratic program with rational data is solvable exactly in
polynomial time, and its optimal value is a rational of polynomial size
\cite{KozlovTarasovKhachiyan1980}. If a certificate is supplied, the value is
instead $q_{P_r}(v)$ for a leaf containing $v$, found by descending the
partition tree. In the message recurrence and in the calibrated tables, every
entry equals the value of a sum of factor values and unary corrections at one
assignment, with each factor and each correction appearing at most once. A
sum of $m_0$ rationals with numerators and denominators at most $2^{\beta}$
has denominator at most $2^{m_0\beta}$ and numerator at most
$m_0 2^{m_0\beta}$. Comparisons and additions on such numbers are polynomial in
their length. No common denominator across different active sets is needed.
\end{proof}

\begin{lemma}[Height of a unique optimizer over a polytope]\label{lem:cr-height}
Let $P=\{x\in\R^n:Gx\le\gamma\}$ be a nonempty bounded polyhedron with
integral $G,\gamma$, and let $f(x)=(\frac12x^{\mathsf T}Hx+h^{\mathsf T}x)/D_f$ with
integral symmetric $H$, integral $h$ and a positive integer $D_f$. Put
$\beta=\max\{1,\max|H_{ij}|,\max|G_{kj}|\}$ and
$\Lambda_0=(2n\beta^2)^n$. If $f$ has a unique minimizer $x^*$ on $P$, then
$x^*=w/q$ with $w\in\Z^n$ and an integer $1\le q\le\Lambda_0$, and $f(x^*)$ has
denominator at most $2D_f\Lambda_0^2$.
\end{lemma}

\begin{proof}
Let $\mathcal E$ be the rows active at $x^*$ and
$\mathcal T=\{d:G_{\mathcal E}d=0\}$. For $d\in\mathcal T$ and small
$\epsilon>0$, both $x^*\pm\epsilon d$ lie in $P$. Optimality gives
$\nabla f(x^*)^{\mathsf T}d=0$ and $d^{\mathsf T}Hd\ge0$. If
$d^{\mathsf T}Hd=0$ for some $d\in\mathcal T\setminus\{0\}$, then
$f(x^*+\epsilon d)=f(x^*)$, contradicting uniqueness. So $H$ is positive
definite on $\mathcal T$. Choose a maximal linearly independent set
$\mathcal E'\subseteq\mathcal E$ of $r\le n$ rows. Since
$\nabla f(x^*)\perp\mathcal T$, there is $\mu$ with
$Hx^*+h+G_{\mathcal E'}^{\mathsf T}\mu=0$, and $G_{\mathcal E'}x^*=\gamma_{\mathcal E'}$.
The saddle matrix
$\left(\begin{smallmatrix}H&G_{\mathcal E'}^{\mathsf T}\\G_{\mathcal E'}&0\end{smallmatrix}\right)$
is nonsingular: if $Hx+G_{\mathcal E'}^{\mathsf T}\mu=0$ and $G_{\mathcal E'}x=0$,
then $x\in\mathcal T$ and $x^{\mathsf T}Hx=0$, so $x=0$ and then $\mu=0$. Its
entries are integers of absolute value at most $\beta$ and its order is at most
$2n$, so Hadamard's inequality bounds its determinant by
$(\sqrt{2n}\beta)^{2n}=\Lambda_0$. Cramer's rule gives the common denominator
$q$, and $f(x^*)=(w^{\mathsf T}Hw+2qh^{\mathsf T}w)/(2q^2D_f)$.
\end{proof}

\begin{theorem}[Certified convex recourse]\label{thm:cr}
In the model of Definition~\ref{def:cr-model}, let each block carry a
certificate on the continuous hull of its attachment box, or none. Put
$\sigma_t=0$ for a block without certificate and otherwise let $\sigma_t$ be
the vector of Proposition~\ref{prop:vf-curv}(b). Define
\begin{equation}\label{eq:cr-L}
L_i=(H_0)_{ii}-\sum_{t:\,i\in S_t}\sigma_{t,j(t,i)},\qquad L=\max_iL_i,
\end{equation}
and let $S$ be the total encoding length of data, decomposition and
certificates.
\begin{enumerate}
\item[(i)] For each $i$, $s\mapsto V(\ldots,s,\ldots)-L_is^2/2$ is concave on
every coordinate-$i$ segment of $\widehat Z$. All certificates are checked in
$\poly(S)$ time.
\item[(ii)] If $L\le0$, endpoint dynamic programming with two labels per
coordinate returns an exact optimizer of $F$ in $2^{O(p)}\poly(S)$ bit
operations. No growth or uniqueness is needed.
\item[(iii)] Suppose $L>0$ and $F$ has quadratic growth $g>0$ in all original
variables. With $\kappa=\max\{1,L/g\}$, the corrected-grid algorithm applied to
$V$ returns a feasible original point with a certified gap at most $2^{-q}$ in
$f(p,\kappa)(S+q+1)^C$ bit operations, and an exact rational optimizer in
$f_1(p,\kappa)(S+1)^{C_1}$ bit operations, with absolute $C,C_1$. Neither $g$
nor $\kappa$ is an input, and every certified lower bound is valid without
growth or uniqueness.
\item[(iv)] $L_i\le(H_0)_{ii}$ for every $i$. Thus the parameter never exceeds
the one obtained by treating all variables as retained.
\end{enumerate}
\end{theorem}

\begin{proof}
(i) The segment $z+se_i$ maps to the coordinate-$j(t,i)$ segment through
$z_{S_t}$ for $t$ with $i\in S_t$, and leaves the other factors constant.
Along it, $q_0$ has second derivative $(H_0)_{ii}$, and
$\phi_t+\sigma_{t,j(t,i)}s^2/2$ is concave by Proposition~\ref{prop:vf-curv}.
A sum of concave functions is concave. The checking cost was discussed after
Definition~\ref{def:leaf}.

(ii) If $L\le0$, $V$ is concave along every coordinate segment of
$\widehat Z$. Replacing coordinates one at a time by a better endpoint does not
increase $V$, so a vertex of $Z$ (integral after inward rounding) is optimal.
Tree dynamic programming over the two endpoint labels, with factor values from
Lemma~\ref{lem:cr-bits}, finds it. Private coordinates are recovered by one
convex quadratic program per block. Every bag table has at most $2^p$ entries.

(iii) The proofs of Lemmas~\ref{lem:round}, \ref{lem:contraction},
\ref{lem:localize}, \ref{lem:states}, of Proposition~\ref{prop:filter}, and of
the capped schedule use the objective only through upper coordinate
curvature on the continuous hull, exact finite factor tables and, for the
quantitative lemmas, global quadratic growth. By (i), $V$ has upper
coordinate curvature $L$. Each $\phi_t$ is assigned to a bag containing $S_t$.
By Lemma~\ref{lem:cr-growth}, $V$ has growth $g$ with the unique minimizer
$z^*$. Hence all those statements hold for $V$ with $\kappa=\max\{1,L/g\}$.
Every incumbent is a retained grid point $z$; its lift $(z,y(z))$ is feasible
and has value $V(z)$. In the arithmetic part of Theorem~\ref{thm:approx},
grid coordinates at stage $j$ of trial $\mu$ have bit length polynomial in
$S+j+\mu K$, and Lemma~\ref{lem:cr-bits} bounds all table entries. There are at
most $T K^p$ value-factor evaluations per stage. The bound $f(p,\kappa)(S+q+1)^C$
follows exactly as in Theorem~\ref{thm:approx}.

For exact output, fix the integer retained coordinates of $z^*$. The
remaining problem is a quadratic over a bounded rational polytope, the product
of the continuous retained box and the sets $Y_t$, and has the unique
minimizer given by the continuous part of $(z^*,y^*)$. Scale its constraint
rows and objective to integers. The numbers $\beta$ and $D_f$ of
Lemma~\ref{lem:cr-height} depend only on the continuous Hessian block, the
constraint matrix and the objective denominator, not on the fixed integer
values. Lemma~\ref{lem:cr-height} therefore gives computable integers $R_0$
and $V_0$ of polynomial bit length such that every coordinate of $z^*$ has
denominator at most $R_0$ and $\OPT$ has denominator at most $V_0$. Run the
approximation with $\epsilon=2^{-q}$ for $q=1,2,4,\ldots$. After each run, if
$U-b<1/V_0^2$, the interval $[b,U]$ contains at most one rational of
denominator at most $V_0$, because two such rationals differ by at least
$1/V_0^2$; it contains $\OPT$, and continued fractions find it in polynomial
time. Around each coordinate of the incumbent take the window of radius
$1/(4R_0^2)$; it contains at most one rational of denominator at most $R_0$. If
all coordinates reconstruct to a point $\tilde z\in Z$, solve the private
programs at $\tilde z$ and accept if the lifted value equals the isolated
$\OPT$. An accepted point is feasible with optimal value, hence optimal,
independently of growth. Once
$\epsilon\le\min\{1/(4V_0^2),\,g/(32R_0^4)\}$, growth puts the incumbent within
$1/(\sqrt{32}R_0^2)<1/(4R_0^2)$ of $z^*$ in every coordinate, so acceptance
occurs. Since $1/g\le\kappa/L$ and $\log(1/L)\le\poly(S)$, this happens with
$q\le\poly(S)+O(\log\kappa)$, and doubling overshoots by a factor at most two.

(iv) $\sigma_t\ge0$ by Proposition~\ref{prop:vf-curv}. By
Lemma~\ref{lem:cr-growth} the growth of $V$ is at least $g$.
\end{proof}

\begin{remark}[Three instances]\label{rem:cr-instances}
With no certificates, $L_i=(H_0)_{ii}$: changing active sets are handled
exactly, but direct stiffness is not cancelled. For $M(y-z)^2$ with
$y\in[0,1]$ and $z\in[0,1]$, $(H_0)_{zz}=2M$. A certificate with one leaf is
a globally affine response. Then $\phi_t=q_t$ is an explicit quadratic with
Hessian $-B_t^{\mathsf T}C_tB_t$, which can be merged into $q_0$; the reduced
problem is again an explicit rational box quadratic. For $M(y-z)^2$ the
response $y=z$ gives $\sigma=2M$ and $L_z=0$. A certificate with several
leaves is a piecewise-affine response. It certifies
$\sigma_j=\min_r(B_r^{\mathsf T}CB_r)_{jj}$, the exact value by
Proposition~\ref{prop:vf-curv}(c). Blocks of different kinds can be mixed. The
local partitions are never overlaid; their product can have exponentially many
cells. Parametric-QP critical regions are classical
\cite{BemporadEtAl2002}; the contribution here is their curvature accounting
inside the conditioned sparse algorithm.
\end{remark}

\begin{corollary}[Affine condensation]\label{cor:cr-affine}
Suppose every block has a one-leaf certificate on its attachment box and the
retained coordinates are continuous. Let
$M=I+\sum_tE_t^{\mathsf T}B_t^{\mathsf T}B_tE_t$ and let
$H_{\mathrm{red}}=H_0-\sum_tE_t^{\mathsf T}B_t^{\mathsf T}C_tB_tE_t$ be the reduced
Hessian. Let $\nu=\max\{0,-\lambda_{\min}(\nabla^2F)\}$.
\begin{enumerate}
\item[(a)] $V(z)-\OPT\ge g(z-z^*)^{\mathsf T}M(z-z^*)$ and
$H_{\mathrm{red}}\succeq-\nu M$.
\item[(b)] If $D$ is a positive diagonal matrix with $M\succeq D$, choose dyadic
$r_i>0$ with $1\le r_i^2D_{ii}<4$ and substitute $z=\operatorname{diag}(r)w$.
The new box is rational, the graph is unchanged, the growth constant in $w$ is
at least $g$, and the upper coordinate curvature in $w$ is at most
$4\max_i(H_{\mathrm{red}})_{ii}/D_{ii}$.
\end{enumerate}
\end{corollary}

\begin{proof}
(a) The first inequality is Lemma~\ref{lem:cr-growth}. With $J$ the derivative
of $z\mapsto(z,\bar y(z))$, $H_{\mathrm{red}}=J^{\mathsf T}\nabla^2F\,J$ and
$J^{\mathsf T}J=M$, so $H_{\mathrm{red}}\succeq-\nu J^{\mathsf T}J$. (b) Growth in
$w$ is at least
$g\sum_ir_i^2D_{ii}(w_i-w_i^*)^2\ge g\norm{w-w^*}^2$, and the diagonal entry
in $w$ is $r_i^2(H_{\mathrm{red}})_{ii}<4(H_{\mathrm{red}})_{ii}/D_{ii}$ when it is
positive.
\end{proof}

Rescaling an integer coordinate changes its lattice, so part~(b) is stated for
continuous coordinates. The inequality in (a) bounds negative curvature only;
it gives no upper bound on $(H_{\mathrm{red}})_{ii}$.

\begin{corollary}[A negative-curvature regime]\label{cor:cr-nu}
In Theorem~\ref{thm:cr}(iii), if the full Hessian of $F$ has
$\nu=\max\{0,-\lambda_{\min}\}>0$ and $L\le C_0\nu$, then
$\kappa\le\max\{1,C_0\}\max\{1,\nu/g\}$. Membership in this regime can be
checked: a rational $\beta$ with $\nu\le\beta<2\nu$ is computable in polynomial
time.
\end{corollary}

\begin{proof}
The bound on $\kappa$ is immediate. For $\beta$, write $H=\nabla^2F$ and let
$D_0H$ be integral. First test $H\succeq0$ exactly; if it fails, $\nu>0$.
Start from $\beta_0=\max_i\sum_j|H_{ij}|\ge\norm H\ge\nu$ and halve while
$H+(\beta/2)I\succeq0$. The final $\beta$ satisfies $H+\beta I\succeq0$ and
$H+(\beta/2)I\not\succeq0$, that is $\nu\le\beta<2\nu$. If $d=\operatorname{rank}H$,
the product of the nonzero eigenvalues is, up to sign, a coefficient of the
characteristic polynomial with denominator dividing $D_0^d$, so its absolute
value is at least $D_0^{-d}$. Each nonzero eigenvalue is at most $\beta_0$ in
absolute value, so $\nu\ge D_0^{-d}\beta_0^{1-d}$, and the number of halvings
is $\poly(S)$. Each step is an exact symmetric elimination.
\end{proof}

The algorithm itself never uses $\nu$ or $\beta$; it computes $L$. The
hypothesis $L\le C_0\nu$ is substantive: Proposition~\ref{prop:cr-limit} below
shows that it can fail for every admissible choice of private blocks.

\subsection{Recognizing a globally affine response}

The affine instance needs no supplied map. Let $Y=[l,u]$ with $l<u$, let
$\Pi=\{v:Wv\le\omega\}$ be a rational polytope with nonempty interior and a
rational interior point $v_0$ (the midpoint of an attachment box), and let
$f(y,v)=\frac12y^{\mathsf T}Cy+y^{\mathsf T}Kv+c^{\mathsf T}y$ with $C\succeq0$.
An \emph{affine optimal selector} is an affine map $\bar y$ with
$\bar y(v)\in\arg\min_{Y}f(\cdot,v)$ for every $v\in\Pi$.

\begin{theorem}[Recognition of affine selectors]\label{thm:cr-recog}
One exact convex quadratic program and one linear program decide whether an
affine optimal selector exists. If one exists, they return a rational selector
of polynomial size together with affine multipliers forming a one-leaf
certificate on $\Pi$. Singular $C$ and nonunique responses are allowed.
\end{theorem}

\begin{proof}
\emph{Common central gradient.} Let $\hat y$ be an optimizer at $v_0$ and
$h_0=Kv_0+c$. For any optimizer $y'$ at $v_0$, the function
$\tau\mapsto f(\hat y+\tau(y'-\hat y),v_0)$ is a convex quadratic that is
constant on $[0,1]$, because the optimal set is convex. Its quadratic
coefficient $(y'-\hat y)^{\mathsf T}C(y'-\hat y)$ vanishes, so $C(y'-\hat y)=0$
since $C\succeq0$. All central optimizers therefore have the same gradient
$g_0=C\hat y+h_0$; this is a classical property of convex programs
\cite{Mangasarian1988}. Let $K_l=\{i:g_{0i}>0\}$, $K_u=\{i:g_{0i}<0\}$,
$J=\{i:g_{0i}=0\}$. By the box KKT conditions, every central optimizer has
$y_i=l_i$ for $i\in K_l$ and $y_i=u_i$ for $i\in K_u$.

\emph{Necessity.} We use one fact: if an affine function is nonnegative on
$\Pi$ and vanishes at an interior point, it vanishes identically. Let
$\bar y$ be an affine optimal selector and write
$\Gamma(v)=C\bar y(v)+Kv+c$, an affine map. Since $\bar y(v_0)$ is a central
optimizer, $\Gamma(v_0)=g_0$. For $i\in K_l$, $\bar y_i-l_i\ge0$ on $\Pi$ and
vanishes at $v_0$, so $\bar y_i\equiv l_i$; optimality at every $v$ then needs
$\Gamma_i\ge0$ on $\Pi$. Symmetrically, $\bar y_i\equiv u_i$ and $\Gamma_i\le0$
for $i\in K_u$. Let $j\in J$, so $\Gamma_j(v_0)=0$. If $\bar y_j\equiv l_j$,
optimality needs $\Gamma_j\ge0$ on $\Pi$, so $\Gamma_j\equiv0$; the case
$\bar y_j\equiv u_j$ is symmetric. Otherwise $\bar y_j$ is not identically at a
bound, so by the fact above it attains no bound at an interior point. Then
optimality gives $\Gamma_j=0$ on the interior of $\Pi$, hence
$\Gamma_j\equiv0$. Thus every affine optimal selector satisfies
\begin{equation}\label{eq:recog-system}
\begin{array}{ll}
\bar y_i\equiv l_i,\ \Gamma_i\ge0\text{ on }\Pi & (i\in K_l),\\
\bar y_i\equiv u_i,\ \Gamma_i\le0\text{ on }\Pi & (i\in K_u),\\
\Gamma_j\equiv0,\ l_j\le\bar y_j\le u_j\text{ on }\Pi & (j\in J).
\end{array}
\end{equation}

\emph{Sufficiency.} If an affine $\bar y$ satisfies~\eqref{eq:recog-system},
define lower multipliers $\Gamma_i$ on $K_l$ and upper multipliers $-\Gamma_i$ on
$K_u$, and zero otherwise. These are affine and nonnegative on $\Pi$, and
satisfy stationarity and complementarity identically. This is a certificate
on $\Pi$; by Lemma~\ref{lem:leaf}(a), $\bar y(v)$ is optimal for every
$v\in\Pi$.

\emph{Linear program.} Write $\bar y(v)=a+Bv$. Then
$\Gamma(v)=(Ca+c)+(CB+K)v$ has coefficients linear in $(a,B)$. Identities in
\eqref{eq:recog-system} are linear equations. Each condition ``an affine
function is nonnegative on $\Pi$'' is replaced by its Farkas certificate
$\pi\ge0$, $W^{\mathsf T}\pi=-(\text{linear part})$,
$\omega^{\mathsf T}\pi\le(\text{constant part})$, which is linear in $(a,B,\pi)$.
For a box $\Pi=v_0+\prod_j[-\rho_j,\rho_j]$, auxiliary variables bounding
$|B_{ij}|$ and $|(CB+K)_{ij}|$ give the exact ranges instead. The resulting
system has polynomial size. It is feasible exactly when an affine optimal
selector exists. Exact convex quadratic programming
\cite{KozlovTarasovKhachiyan1980} and exact linear programming
\cite{GrotschelLovaszSchrijver1988} run in polynomial time and return rational
solutions of polynomial size. If $k=0$, any central optimizer is a constant
selector.
\end{proof}

Fixing the active bounds of an arbitrary central optimizer is not enough when
$C$ is singular. For $(y_1+y_2-z)^2$ with $y\in[0,1]^2$ and $z\in[0,2]$, the
central optimizer $(0,1)$ would fix both coordinates, yet the common central
gradient is zero and the theorem returns the selector $(z/2,z/2)$. For
$(y-z)^2$ with $y\in[0,1]$ and $z\in[-1,2]$, the clipped response is not
affine, and the linear program is infeasible. A negative answer rejects the
affine class for this block, not the optimization problem; the block may
still enter Theorem~\ref{thm:cr} with a piecewise certificate or with
$\sigma=0$. Choosing a useful partition into private blocks is not addressed.

\subsection{Examples}

\begin{example}[Three pieces, cancellation $2M$]\label{ex:cr-fm}
For $M\ge1$ let $z\in[0,3/4]$, $y\in[0,1]^2$, and
$f_M=M(y_1+y_2-z)^2+(y_1-2z+\frac12)^2$. The retained-only part is
$Mz^2+(2z-\frac12)^2$, so $(H_0)_{zz}=2M+8$, and
$C=\left(\begin{smallmatrix}2M+2&2M\\2M&2M\end{smallmatrix}\right)\succ0$,
$K=(-2M-4,-2M)^{\mathsf T}$, $c=(1,0)^{\mathsf T}$. The unique responses and
values are
\[
\begin{array}{llll}
z\in[0,\tfrac14]:& y=(0,z),& \min_yf_M=(\tfrac12-2z)^2,& B^{\mathsf T}CB=2M,\\
z\in[\tfrac14,\tfrac12]:& y=(2z-\tfrac12,\tfrac12-z),& \min_yf_M=0,& B^{\mathsf T}CB=2M+8,\\
z\in[\tfrac12,\tfrac34]:& y=(\tfrac{(M+2)z-1/2}{M+1},0),& \min_yf_M=\tfrac{M(z-1/2)^2}{M+1},
& B^{\mathsf T}CB=\tfrac{2(M+2)^2}{M+1}.
\end{array}
\]
On the three intervals the private gradients are $(1-4z,0)$, $(0,0)$ and
$(0,2M(z-\frac12)/(M+1))$, all with the signs required at the active lower
bounds; in the third interval $y_1\in[\frac12,1)$. These are certificates. By
Proposition~\ref{prop:vf-curv}, $\sigma=2M$, and the certified curvature is
$2M+8-2M=8$, while the direct bound is $2M+8$. The piece second derivatives
$8,0,2M/(M+1)$ agree with $(H_0)_{zz}-B^{\mathsf T}CB$. The response is unique
and changes slope, so no affine selector exists.
\end{example}

\begin{proposition}[A ladder with growing negative inertia]\label{prop:cr-ladder}
For $m\ge1$, $M\ge1$, $a=\frac15$ and $\eta=\frac1{16}$, let
$z_i\in[0,\frac34]$, $v_i\in[0,1]$, $y_{1i},y_{2i}\in[0,1]$ and
\[
F=\sum_{i=1}^m\Bigl[z_i^2+f_M(y_{1i},y_{2i},z_i)-v_i^2+3v_i-(z_i-a)v_i\Bigr]
+\eta\sum_{i=1}^{m-1}\Bigl[(z_{i+1}-z_i)^2+(v_{i+1}-v_i)^2\Bigr].
\]
Then $F$ has the unique minimizer $z_i=a$, $v_i=0$, $y_{1i}=0$, $y_{2i}=a$, with
$\OPT=m/20$ and growth constant $g=1/12$. Its Hessian has exactly $m$
negative eigenvalues and $2\le\nu\le\sqrt5$. With the private blocks
$(y_{1i},y_{2i})$, the retained bags $\{z_i,v_i,z_{i+1},v_{i+1}\}$ have size at
most four, and the certificates of Example~\ref{ex:cr-fm} give $L\le41/4$. The
direct diagonal at $z_i$ is at least $2M+10$, while $g\le33/16$.
\end{proposition}

\begin{proof}
Put $t=z-a$, $u=y_1$, $w=y_2-a$, $r=u+w-t$ and $s=u-2t$ for one index. Direct
expansion gives
\[
z^2+f_M-\tfrac1{20}=t^2+Mr^2+s^2+\tfrac u5.
\]
Also $-v^2+3v-tv\ge v^2$, because $v(3-2v-t)\ge0$ for $v\in[0,1]$ and
$t\le\frac{11}{20}$. Since $u\ge0$ and the path terms are nonnegative,
\[
F-\tfrac m{20}\ge\sum_i\bigl(t_i^2+Mr_i^2+s_i^2+v_i^2\bigr).
\]
From $u=s+2t$ and $w=r-s-t$, $u^2\le2s^2+8t^2$ and $w^2\le3(r^2+s^2+t^2)$, so
$t^2+u^2+w^2\le12(t^2+r^2+s^2)\le12(t^2+Mr^2+s^2)$. Hence
$F-\OPT\ge\frac1{12}\norm{x-x^*}^2$, and equality $F=\OPT$ forces
$t=r=s=v=0$, $u=w=0$.

The terms $z_i^2-v_i^2-(z_i-a)v_i$ contribute the block
$\left(\begin{smallmatrix}2&-1\\-1&-2\end{smallmatrix}\right)$ on each pair
$(z_i,v_i)$, with eigenvalues $\pm\sqrt5$. All other quadratic terms, namely
the squares in $f_M$ and the path terms, are positive semidefinite. So
$\lambda_{\min}\ge-\sqrt5$. The direction $v=m^{-1/2}(1,\ldots,1)$, all
other coordinates zero, has Rayleigh quotient $-2$, so $\nu\ge2$. On the
$m$-dimensional $v$-subspace the form is $-2I+2\eta L_{\mathrm{path}}$, and
$2\eta\norm{L_{\mathrm{path}}}\le\frac12$, so it is negative definite; on the
subspace $\{v=0\}$, of codimension $m$, the form is a sum of squares. Hence the
negative inertia is exactly $m$.

The retained diagonal at $z_i$ is
$(H_0)_{z_iz_i}=2+(2M+8)+2\eta\deg_i$ with $\deg_i\le2$, and at $v_i$ it is
$-2+2\eta\deg_i<0$. Each block has $\sigma=2M$ by Example~\ref{ex:cr-fm}, so
$L_{z_i}\le10+4\eta=41/4$. Raising $v_1$ from $0$ to $1$ at the minimizer
increases $F$ by at most $2+\eta$, which gives $g\le33/16$.
\end{proof}

Thus the certified ratio is $L/g\le123$ for all $m,M$, whereas the original
coordinate ratio is at least $(2M+10)\cdot16/33$. Every instance also has the
explicit nonnegative expansion above; the family separates parameters and is
not offered as a hard benchmark.

\begin{remark}[A one-leaf ladder]\label{rem:cr-oneleaf}
With all variables in $[0,1]$, $a=\frac13$, $\eta=\frac1{16}$ and $M_0\ge1$, let
\[
F=\sum_{i=1}^m\Bigl[(u_i-a)^2-(u_i-a)v_i-v_i^2+3v_i+M_0(y_i-u_i)^2\Bigr]
+\eta\sum_{i=1}^{m-1}\Bigl[(u_{i+1}-u_i)^2+(v_{i+1}-v_i)^2\Bigr].
\]
The response $y_i=u_i$ with zero multipliers is a one-leaf certificate, so
$\sigma=2M_0$ cancels the direct term $2M_0$, and $L\le2+4\eta=\frac94$. Since
$-(u-a)v-v^2+3v\ge v^2$ for $u-a\le\frac23$, we get
$F\ge\norm{u-a\mathbf 1}^2+\norm v^2+M_0\norm{y-u}^2$, and with
$\norm{y-a\mathbf 1}^2\le2\norm{y-u}^2+2\norm{u-a\mathbf 1}^2$ this gives growth
$\frac13$ at the unique minimizer $u=y=a\mathbf 1$, $v=0$. As in
Proposition~\ref{prop:cr-ladder}, the form is negative definite on the
$v$-subspace and positive semidefinite on $\{v=0\}$, so the negative inertia is
exactly $m$, and $2\le\nu\le\sqrt5$. The direct diagonal at $u_i$ is at least
$2M_0+2$.
\end{remark}

\subsection{Limits of global recourse curvature}

Call a retained set $\mathcal R\subseteq\{1,\ldots,n\}$ \emph{admissible} if the
principal Hessian block of the remaining coordinates is positive
semidefinite, so that they form convex private blocks. For admissible
$\mathcal R$ the reduced function $V_{\mathcal R}$ is intrinsic: it does not
depend on how the private part is grouped or certified. Let
$L_{\mathcal R}$ be the least $L$ for which~\eqref{eq:semiconcavity} holds for
$V_{\mathcal R}$ in every retained coordinate, and $g_{\mathcal R}$ the largest
growth constant of $V_{\mathcal R}$. By Proposition~\ref{prop:vf-curv}(c), no
certificate yields a parameter below $L_{\mathcal R}/g_{\mathcal R}$.

\begin{proposition}[Recourse cannot replace $L/g$ by $\nu/g$ in general]\label{prop:cr-limit}
For $M\ge1$ let $x\in[0,2]$, $y\in[1,3]$ and
\[
F_M(x,y)=M(x-y)^2-\tfrac18(x+y-2)^2+2(y-1).
\]
Then $(1,1)$ is the unique minimizer, $\OPT=0$,
$F_M(x,y)-\OPT\ge\frac18\bigl((x-1)^2+(y-1)^2\bigr)$, and $\nu=\frac12$; thus
$\nu/g\le4$ with $g=\frac18$. Nevertheless, for every
admissible retained set $\mathcal R$ and every positive diagonal rescaling of
the retained coordinates,
\[
L_{\mathcal R}/g_{\mathcal R}\ \ge\ (4M-\tfrac12)/3\ \ge\ M.
\]
\end{proposition}

\begin{proof}
Put $\alpha=x-1\in[-1,1]$ and $\beta=y-1\in[0,2]$, so
$F_M=M(\alpha-\beta)^2-\frac18(\alpha+\beta)^2+2\beta$. Using $M\ge1$ and
$2\beta\ge\beta^2$ on $[0,2]$,
\[
F_M-\tfrac18(\alpha^2+\beta^2)\ge\tfrac34\alpha^2-\tfrac94\alpha\beta+\tfrac74\beta^2
=\tfrac34\bigl(\alpha-\tfrac32\beta\bigr)^2+\tfrac1{16}\beta^2\ge0.
\]
This proves growth $\frac18$, uniqueness and $\OPT=0$. The Hessian
$\left(\begin{smallmatrix}2M-1/4&-2M-1/4\\-2M-1/4&2M-1/4\end{smallmatrix}\right)$
has eigenvalues $4M$ on $(1,-1)$ and $-\frac12$ on $(1,1)$; so $\nu=\frac12$,
and $\{x,y\}$ is not a convex block. The admissible sets are $\{x,y\}$,
$\{y\}$ and $\{x\}$. Note $F_M(2,2)=\frac32$.

$\mathcal R=\{x,y\}$: $V=F_M$ has $L_{\mathcal R}=2M-\frac14$. Under
$x=r_1x'$, $y=r_2y'$ the coordinate curvatures become $r_i^2(2M-\frac14)$, and
the point $(2,2)$ shows $g'\le\frac32/(r_1^{-2}+r_2^{-2})\le\frac32\min_ir_i^2$.

$\mathcal R=\{y\}$: $F_M$ is strictly convex in $x$, and
$\partial_xF_M(2,y)=2M(2-y)-y/4\le0$ for $y\ge y_c=4M/(2M+\frac14)$, where
$1<y_c<2$. So on $[y_c,3]$ the response is $x=2$ and
$V(y)=F_M(2,y)$ has second derivative $2M-\frac14$. Hence
$L_{\mathcal R}\ge2M-\frac14$, while $V(2)\le\frac32$ and $V(1)=0$ give
$g_{\mathcal R}\le\frac32$.

$\mathcal R=\{x\}$: $F_M$ is strictly convex in $y$, and
$\partial_yF_M(x,1)=2-(2M+\frac14)(x-1)\ge0$ for
$x\le x_c=1+2/(2M+\frac14)$, where $1<x_c\le2$. So on $[0,x_c]$ the response
is $y=1$ and $V(x)=(M-\frac18)(x-1)^2$ has second derivative $2M-\frac14$.
Again $V(2)\le\frac32$ gives $g_{\mathcal R}\le\frac32$.

In one dimension a rescaling multiplies curvature and growth by the same
factor. In all cases $L_{\mathcal R}/g_{\mathcal R}\ge(2M-\frac14)/\frac32$.
\end{proof}

Taking disjoint copies gives arbitrarily many variables with the same
constants and bag size two. The instance is trivial to solve; the proposition
limits the parameter of Theorem~\ref{thm:cr}, not the problem. The obstruction
is the clipped piece: the private coordinate that cancels the stiffness hits a
bound on a full-dimensional set of parameters. The piece need not be large.

\begin{remark}[The certified cancellation is unstable]\label{rem:cr-unstable}
For $f_\epsilon=M(y-u)^2+\epsilon y$ with $y,u\in[0,1]$, the response is
$y=\max\{0,u-\epsilon/(2M)\}$. For $\epsilon=0$ it is $y=u$, so $\sigma=2M$.
For every $\epsilon>0$ the piece $[0,\epsilon/(2M)]$ has response $y=0$, so
$\sigma=0$ by Proposition~\ref{prop:vf-curv}(c), and the value has second
derivative $2M$ there. The full Hessian, and hence $\nu$, does not depend on
$\epsilon$. Adding $\epsilon\sum_iy_i$ to the ladder of
Remark~\ref{rem:cr-oneleaf} therefore raises its certified curvature from
$\frac94$ to $2M_0+\frac94$. The same happens near $z=0$ in
Example~\ref{ex:cr-fm} if $\epsilon y_2$ is added: on a short interval both
private coordinates sit at their lower bounds, and $\sigma=0$. The certified
quantity is a maximum over pieces, however small the piece.
\end{remark}

The two nontrivial choices fail for different reasons. For
$\mathcal R=\{x\}$ the stiff piece $[0,x_c]$ contains the minimizer, but there
the growth is stiff as well, $V(x)=(M-\frac18)(x-1)^2$; on $[x_c,2]$ the
response is interior and $V$ is concave. So on each piece the ratio of
curvature to growth about $x^*$ is $O(1)$. For $\mathcal R=\{y\}$ the stiff
piece $[y_c,3]$ lies where $V\ge V(y_c)=\frac32-O(1/M)$, so it contains no
point within objective gap about $\frac32$ of the optimum. In both cases the
large constant comes from a region that a localized analysis could treat
separately. A bound in terms of $p$ and $\nu/g$ along this route would need
curvature measured only where near-optimal points can lie. The corrections of
Lemma~\ref{lem:round} may already use the curvature of the current box, but
the first stages still see the stiff pieces, and no such analysis is given
here.

\subsection{The negative-curvature target}

Let $X$ be a continuous box, $F(x)=\frac12x^{\mathsf T}Hx+b^{\mathsf T}x+c_0$, and
$\nu=\max\{0,-\lambda_{\min}(H)\}$. The target is a parameter
$\max\{1,\nu/g\}$ in place of $\kappa$, letting large positive curvature affect
only the encoding length. The target is continuous: with integer variables,
$\nu=0$ still contains convex integer quadratic programming.

\begin{lemma}[Convex energy]\label{lem:cr-energy}
Let $x^*$ minimize $F$ on $X$ with $F(x)-\OPT\ge g\norm{x-x^*}^2$. For every
$c\ge0$ and $x\in X$, with $d=x-x^*$,
\[
F(x)-\OPT\ge\frac{g}{2g+c}\,d^{\mathsf T}(H+cI)d.
\]
For $c=2\nu$ the matrix $H+2\nu I\succeq\nu I$ has $H\preceq H+2\nu I$, so in its
metric $F$ has upper curvature $1$ and growth $g/(2g+2\nu)$.
\end{lemma}

\begin{proof}
First-order optimality on the convex set $X$ gives
$\nabla F(x^*)^{\mathsf T}d\ge0$, so the exact expansion yields
$F(x)-\OPT\ge\frac12d^{\mathsf T}Hd$. Combine it with growth using weights
$2g/(2g+c)$ and $c/(2g+c)$.
\end{proof}

\begin{lemma}[Proximal envelope]\label{lem:cr-envelope}
For $\lambda>0$ and $w\in\R^n$ let
$E_\lambda(w)=\min_{x\in X}\{F(x)+\frac\lambda2\norm{x-w}^2\}$. Then
$E_\lambda-\frac\lambda2\norm\cdot^2$ is concave, so $E_\lambda$ has upper
curvature $\lambda$ in every direction. Under the growth hypothesis of
Lemma~\ref{lem:cr-energy}, $E_\lambda(x^*)=\OPT$ and
\[
E_\lambda(w)-\OPT\ge\frac{g\lambda}{2g+\lambda}\norm{w-x^*}^2.
\]
If $\lambda>\nu$, every evaluation is a strongly convex quadratic program.
\end{lemma}

\begin{proof}
$E_\lambda(w)-\frac\lambda2\norm w^2=\min_x\{F(x)+\frac\lambda2\norm x^2-\lambda x^{\mathsf T}w\}$
is a minimum of affine functions of $w$. Next,
$E_\lambda(w)-\OPT\ge\min_{x\in\R^n}\{g\norm{x-x^*}^2+\frac\lambda2\norm{x-w}^2\}
=\frac{g\lambda}{2g+\lambda}\norm{w-x^*}^2$. Finally $\OPT\le E_\lambda(x^*)\le F(x^*)$.
The last claim holds because $H+\lambda I\succ0$.
\end{proof}

With $\lambda=2\nu$ the envelope has ratio $2+2\nu/g$. Neither lemma gives a
sparse algorithm. Whitening by $H+2\nu I$ destroys the product box, and
independent rounding does not preserve the cancellations of the energy. On a
fixed active set without bounds, the envelope Hessian is
$\lambda H(H+\lambda I)^{-1}$, which is generally dense even for tridiagonal
$H$.

\begin{proposition}[Matched separator moments do not control the error]\label{prop:cr-moment}
For rational $0<h\le\frac12$ let $(s,u_1,u_2,v)\in[0,1]^4$ and
\[
F_h=\underbrace{8\bigl(\tfrac sh-\tfrac{u_1+u_2}2\bigr)^2+\textstyle\sum_{i=1}^2u_i(1-u_i)
+\tfrac{u_1+u_2}{16}}_{f_1(s,u_1,u_2)}
+\underbrace{8\bigl(\tfrac sh-\tfrac14-\tfrac v2\bigr)^2+v(1-v)+\tfrac v{16}}_{f_2(s,v)}.
\]
Then $F_h$ has the unique minimizer $(h/8,0,0,0)$, $\OPT=\frac14$, growth
constant $\frac1{22}$, and $\nu=2$. Let $\mu_1$ put masses
$\frac18,\frac38,\frac38,\frac18$ at
$(s,u_1,u_2)=(0,0,0),(\frac h2,1,0),(\frac h2,0,1),(h,1,1)$, and let $\mu_2$ put
masses $\frac12,\frac12$ at $(s,v)=(\frac h4,0),(\frac{3h}4,1)$. The two
laws of $s$ have equal moments of orders $1,2,3$, all atoms have $s\in[0,h]$
and binary controls, and
\[
\E_{\mu_1}f_1+\E_{\mu_2}f_2=\tfrac3{32}=\OPT-\tfrac5{32}.
\]
Moreover a single positive semidefinite covariance matrix has both bag
covariances as principal submatrices.
\end{proposition}

\begin{proof}
With $\delta=s/h-\frac18-(u_1+u_2+v)/4$, direct expansion gives
\[
F_h-\tfrac14=16\delta^2+2\bigl[(1-v)u_1u_2+v(1-u_1)(1-u_2)\bigr]+\tfrac{u_1+u_2+v}{16},
\]
which is nonnegative on the box and vanishes only at $\delta=0$, $u=v=0$.
Write $\omega=(u_1,u_2,v)$. Since each coordinate is in $[0,1]$,
$F_h-\frac14\ge16\delta^2+\norm\omega^2/16$. Also
$s-h/8=h(\delta+(u_1+u_2+v)/4)$, so with $h\le\frac12$,
$(s-h/8)^2\le\frac12\delta^2+\frac3{32}\norm\omega^2$ and
$\norm{x-x^*}^2\le\frac12\delta^2+\frac{35}{32}\norm\omega^2\le22(16\delta^2+\norm\omega^2/16)$.
The Hessian is a positive semidefinite sum of squares minus $2$ on each
control; the direction $(0,1,-1,0)$ leaves both squares unchanged, so
$\lambda_{\min}=-2$. For the measures, with $\tau=s/h$, both laws of $\tau$ have
moments $\frac12,\frac5{16},\frac7{32}$. At every atom both squares and all
terms $u_i(1-u_i)$, $v(1-v)$ vanish, so the expectations equal
$\E(u_1+u_2)/16+\E v/16=\frac1{16}+\frac1{32}$. For the covariance, take mean
$(h/2,\frac12,\frac12,\frac12)$ and
$\Sigma=\frac1{16}aa^{\mathsf T}+\frac3{16}bb^{\mathsf T}$ with
$a=(h,1,1,2)^{\mathsf T}$, $b=(0,1,-1,0)^{\mathsf T}$; its restrictions to
$\{s,u_1,u_2\}$ and $\{s,v\}$ equal the covariances of $\mu_1$ and $\mu_2$.
\end{proof}

For moments conditioned on one global cell of widths $w_i$, the relaxed value
is at least $F(m)-\frac\nu2\operatorname{tr}\Sigma\ge F(m)-\frac\nu8\sum_iw_i^2$.
Proposition~\ref{prop:cr-moment} shows that bag-local measures matching
separator moments up to order three do not inherit this bound: on mesh $2h$
every atom lies in a cell whose widths tend to zero, yet the gap stays
$\frac5{32}$ at fixed $\nu/g=44$. The witness does not apply to adaptive
filtering or to all moment hierarchies; matching the fourth moment of $s$
already excludes it.
